\documentclass[11pt]{article}
\usepackage{amsmath, amssymb, physics}
\usepackage{geometry}
\geometry{margin=1in}

\title{Step-by-Step Notes Toward Quantum Field Theory}
\author{}
\date{}

\begin{document}

\maketitle

\section{Introduction}
These notes build quantum field theory (QFT) step by step starting from non-relativistic quantum mechanics (QM). We proceed slowly and carefully, emphasizing physical meaning and mathematical structure.

\section{Review of Non-Relativistic Quantum Mechanics}

\subsection{States and Observables}
A quantum system is described by a state vector $\ket{\psi(t)}$ in a Hilbert space.

Observables are represented by Hermitian operators $\hat{A}$.

The expectation value is:
\begin{equation}
\langle A \rangle = \bra{\psi} \hat{A} \ket{\psi}.
\end{equation}

\subsection{Schr\"odinger Equation}
Time evolution is governed by:
\begin{equation}
i \hbar \frac{d}{dt} \ket{\psi(t)} = \hat{H} \ket{\psi(t)}.
\end{equation}

In position space:
\begin{equation}
i \hbar \frac{\partial}{\partial t} \psi(x,t) = \left(-\frac{\hbar^2}{2m} \nabla^2 + V(x)\right) \psi(x,t).
\end{equation}

\subsection{Canonical Commutation Relations}
Position and momentum satisfy:
\begin{equation}
[x_i, p_j] = i \hbar \delta_{ij}.
\end{equation}

\section{Motivation for Quantum Field Theory}

Non-relativistic QM has key limitations:
\begin{itemize}
\item It is not Lorentz invariant.
\item It cannot describe particle creation/annihilation.
\end{itemize}

Relativity requires energy-momentum relation:
\begin{equation}
E^2 = p^2 + m^2.
\end{equation}

Attempting to quantize this directly leads to problems, motivating fields instead of wavefunctions.

\section{From Particles to Fields}

In QFT, the fundamental object is a field $\phi(x,t)$ defined at every point in spacetime.

\subsection{Classical Field Theory}

A classical scalar field has Lagrangian density:
\begin{equation}
\mathcal{L} = \frac{1}{2}(\partial_\mu \phi)(\partial^\mu \phi) - \frac{1}{2} m^2 \phi^2.
\end{equation}

The action is:
\begin{equation}
S = \int d^4x \, \mathcal{L}.
\end{equation}

Using the Euler-Lagrange equation:
\begin{equation}
\frac{\partial \mathcal{L}}{\partial \phi} - \partial_\mu \left(\frac{\partial \mathcal{L}}{\partial (\partial_\mu \phi)}\right) = 0,
\end{equation}

we obtain the Klein-Gordon equation:
\begin{equation}
(\Box + m^2)\phi = 0.
\end{equation}

\section{Canonical Quantization of Fields}

We now promote the field to an operator:
\begin{equation}
\phi(x) \rightarrow \hat{\phi}(x).
\end{equation}

Define conjugate momentum:
\begin{equation}
\pi(x) = \frac{\partial \mathcal{L}}{\partial (\partial_0 \phi)} = \dot{\phi}(x).
\end{equation}

Impose commutation relations:
\begin{equation}
[\hat{\phi}(\mathbf{x},t), \hat{\pi}(\mathbf{y},t)] = i \hbar \delta^{(3)}(\mathbf{x}-\mathbf{y}).
\end{equation}

\section{Mode Expansion}

Solutions can be expanded in Fourier modes:
\begin{equation}
\hat{\phi}(x) = \int \frac{d^3 p}{(2\pi)^3} \frac{1}{\sqrt{2E_p}} \left(a_{\mathbf{p}} e^{-ipx} + a^\dagger_{\mathbf{p}} e^{ipx}\right).
\end{equation}

Operators satisfy:
\begin{equation}
[a_{\mathbf{p}}, a^\dagger_{\mathbf{q}}] = (2\pi)^3 \delta^{(3)}(\mathbf{p}-\mathbf{q}).
\end{equation}

\section{Interpretation: Particles as Quanta}

The vacuum state $\ket{0}$ satisfies:
\begin{equation}
a_{\mathbf{p}} \ket{0} = 0.
\end{equation}

A one-particle state is:
\begin{equation}
\ket{\mathbf{p}} = a^\dagger_{\mathbf{p}} \ket{0}.
\end{equation}

Thus, particles arise as excitations of the field.

\section{Next Steps}

Future topics:
\begin{itemize}
\item Interacting fields
\item Perturbation theory
\item Feynman diagrams
\item Dirac field and spin
\item Gauge theories
\end{itemize}

\end{document}
